\documentclass[11pt,a4paper]{article}
\usepackage[utf8]{inputenc}
\usepackage[margin=1in]{geometry}
\usepackage{amsmath,amssymb,amsfonts,amsthm}
\usepackage{graphicx}
\usepackage{booktabs}
\usepackage{hyperref}
\usepackage{cite}
\usepackage{microtype}

\hypersetup{
    colorlinks=true,
    linkcolor=blue,
    citecolor=blue,
    urlcolor=blue
}

\title{\textbf{Primordial Inflationary Perturbations Without Inflaton Scalar Fields: Conformal Phase Dynamics of the Feynman Path Functional}}

\author{
    \textbf{H. Pavliuk (Hryhorii Pavliuk)}\\
    \textit{Independent Researcher}\\
\small ORCID: \href{https://orcid.org/0009-0000-0899-3463}{0009-0000-0899-3463}\\
    \small Permanent Monograph DOI: \href{https://doi.org/10.5281/zenodo.24226129}{10.5281/zenodo.24226129}\\
    \small Repository: \href{https://github.com/hrpavliuk-byte/Unified-Phase-Dynamics-Feynman-Path-Contours}{Unified-Phase-Dynamics-Feynman-Path-Contours}
}
\date{November 2026}

\begin{document}

\maketitle

\begin{abstract}
We present a first-principles, deductive derivation of the primordial cosmological perturbation spectrum without postulating ad-hoc beyond-the-Standard-Model (BSM) inflaton scalar fields, fine-tuned slow-roll potentials $V(\phi)$, or trans-Planckian initial conditions. In the \textbf{Unified Phase Dynamics of Feynman Path Contours}, the primordial accelerating expansion and its quantum fluctuations originate directly from the non-perturbative conformal phase holonomy of the universal path master-functional $\mathbf{\Psi}[\Sigma]$ over the gauge-spacetime group $\mathcal{G} = U(1)_Y \times SU(2)_L \times SU(3)_c \times \mathrm{Spin}(3,1)$. Quantum metric fluctuations are regularized at the fundamental scale by an entire-order, non-local Schwinger contour form-factor $\mathcal{E}(\Box) = \exp(-s_{\min}^2 \Box^2/\hbar^2)$ ($s_{\min} = 2\ell_{\mathrm{Pl}}$), which eliminates trans-Planckian divergences without introducing Ostrogradsky ghost poles or violating unitarity and reflection positivity (OS2). We show that the conformal mode of the spatial geometry dynamically locks to the master phase current $\mathcal{U}_\mu = \partial_\mu \theta$, generating an effective horizon attractor potential of the Starobinsky class $V(\theta) = \frac{3}{4} M_{\mathrm{Pl}}^2 M^2 \left(1 - e^{-\sqrt{2/3}\theta/M_{\mathrm{Pl}}}\right)^2$ purely from the quantum gravitational trace anomaly and path holonomies. For $N = 56.5$ e-folds of conformal expansion, the theory predicts:
(i) a primordial scalar spectral index $n_s = 1 - \frac{2}{N + 1/2} = \mathbf{0.96491}$, in exact agreement with the Planck 2018 + BAO baseline ($n_s = 0.9649 \pm 0.0042$, $|\Delta n_s| = 0.003\sigma$);
(ii) a tensor-to-scalar ratio $r = \frac{12}{(N + 1/2)^2} = \mathbf{0.00369} = 3.69 \times 10^{-3}$, comfortably satisfying the BICEP/Keck (BK18) + Planck upper bound ($r < 0.036$) while offering a decisive, falsifiable target for next-generation CMB B-mode observatories (LiteBIRD, CMB-S4);
(iii) running of the scalar tilt $\alpha_s = \frac{dn_s}{d\ln k} = - \mathbf{0.00062}$, fully compatible with Planck limits;
(iv) strictly luminal tensor propagation $c_T \equiv c$ with zero isocurvature modes ($\mathcal{S}_{\mathrm{iso}} \equiv 0$) and minimal non-Gaussianity $f_{\mathrm{NL}}^{\mathrm{local}} = \frac{5}{12}(1 - n_s) \approx 0.015$. This work establishes that cosmic inflation is not an arbitrary phenomenological add-on, but an inevitable consequence of non-perturbative quantum phase dynamics on the cosmological horizon.
\end{abstract}

\section{Introduction: The Inflaton Crisis and Conceptual Dilemmas}
Cosmic inflation, originally conceived by Starobinsky \cite{Starobinsky1980}, Guth \cite{Guth1981}, and Linde \cite{Linde1982}, resolves the classical horizon, flatness, and magnetic monopole problems of Big Bang cosmology, while providing a quantum mechanical mechanism for the origin of large-scale structure via vacuum zero-point fluctuations \cite{Mukhanov1981}. Despite this empirical success, conventional inflationary cosmology suffers from profound theoretical difficulties:
\begin{enumerate}
    \item \textbf{The Particle Identity Crisis:} The canonical inflaton field $\phi$ has no counterpart in the $SU(3)_c \times SU(2)_L \times U(1)_Y$ Standard Model of particle physics. Non-minimal Higgs inflation ($\xi H^\dagger H R$) requires an unnaturally large coupling $\xi \sim 10^4$, leading to perturbative unitarity breakdown at the intermediate energy scale $\Lambda_{\mathrm{cut}} \sim M_{\mathrm{Pl}}/\xi \ll M_{\mathrm{Pl}}$ \cite{Burgess2009}.
    \item \textbf{The $\eta$-Problem and UV-Sensitivity:} In supergravity and generic quantum field theories coupled to gravity, radiative quantum corrections generate dimension-6 operators of the form $\Delta V \sim \frac{V_0}{M_{\mathrm{Pl}}^2} \phi^2$, which naturally push the second slow-roll parameter $\eta \equiv M_{\mathrm{Pl}}^2 \frac{V''}{V} \sim \mathcal{O}(1)$. This abruptly halts the accelerated phase after a few e-folds unless fine-tuned to dozens of decimal places \cite{Baumann2009}.
    \item \textbf{The Trans-Planckian Problem:} For an inflationary epoch exceeding $N \approx 60$ e-folds, cosmological modes observed today in the CMB at comoving wavenumber $k \sim 0.05\text{ Mpc}^{-1}$ had physical wavelengths $\lambda_{\mathrm{phys}}(t_i) = \frac{a(t_i)}{k} \ll \ell_{\mathrm{Pl}}$ far below the Planck length at the onset of inflation. Standard QFT on curved space assumes a naive continuum vacuum deep in the trans-Planckian regime \cite{Martin2001}.
    \item \textbf{The Multiverse Deadlock:} High-dimensional scalar potentials $V(\phi, \chi_i)$ inevitably trigger eternal inflation, generating an untestable ``multiverse'' devoid of predictive power.
\end{enumerate}

In this work, we demonstrate that all these dilemmas vanish when cosmic inflation is derived deductively from the \textbf{Unified Phase Dynamics of Feynman Path Contours} \cite{Pavliuk2024_Monograph}. The primordial expansion is driven not by an ad-hoc scalar field with a hand-tuned potential, but by the \textbf{quantum conformal phase holonomy of the universal path functional $\mathbf{\Psi}[\Sigma]$}.

\section{Conformal Phase Dynamics of the Universal Master Functional}

\subsection{The Master Path Functional $\mathbf{\Psi}[\Sigma]$}
In Unified Phase Dynamics, spacetime geometry and matter fields emerge from the universal path master-functional defined over oriented loop space $\mathcal{L}(\mathcal{M})$:
\begin{equation}
\mathbf{\Psi}[\Sigma] = \int \mathcal{D}[A] \mathcal{D}[e] \exp\left( \frac{i}{\hbar} S_{\mathrm{gauge}}[A] + \frac{i}{\hbar} S_{\mathrm{geom}}[e, \omega] \right),
\end{equation}
where $\mathcal{G} = U(1)_Y \times SU(2)_L \times SU(3)_c \times \mathrm{Spin}(3,1)$. The fundamental boundary contour $\Sigma$ is regularized at the minimal quantum area scale $s_{\min} = 2\ell_{\mathrm{Pl}} = 2\sqrt{\hbar G/c^3}$.

The geometry of the early universe is described by the relational Wheeler--DeWitt equation \cite{Pavliuk2024_Monograph, Pavliuk2024_Work13}:
\begin{equation}
\hat{\mathcal{H}} \mathbf{\Psi}[\Sigma] = 0, \quad \hat{\mathcal{H}}_i \mathbf{\Psi}[\Sigma] = 0.
\end{equation}

\subsection{Conformal Metric Decomposition and the Phase Holonomy}
Under the ADM 3+1 decomposition in conformal time $\tau$ ($dt = a(\tau) d\tau$), the primordial metric is written as:
\begin{equation}
g_{\mu\nu} dx^\mu dx^\nu = a^2(\tau) \left[ -(1 + 2\Phi) d\tau^2 + 2 B_i d\tau dx^i + \left( (1 - 2\Psi)\delta_{ij} + 2 E_{ij} \right) dx^i dx^j \right].
\end{equation}
In the absence of anisotropic stress from relativistic shear, $\Phi = \Psi$.

The spatial volume degree of freedom $a(\tau) = e^{\Omega(\tau)}$ couples directly to the collective phase holonomy of the master functional:
\begin{equation}
W_P[\Sigma] = \mathcal{P} \exp\left( i \oint_\Sigma A_\mu dx^\mu \right) = \exp(i \theta(x, \tau)).
\end{equation}
The topological 4-vector phase current is defined as:
\begin{equation}
\mathcal{U}_\mu \equiv \hbar \, \partial_\mu \theta.
\end{equation}

\subsection{Non-Local Schwinger Regularizer and Trace Anomaly}
The quantum path integral over the metric conformal mode $\Omega(x)$ is regularized by the entire non-local Schwinger contour form-factor:
\begin{equation}
\mathcal{E}(\Box) = \exp\left( -\frac{s_{\min}^2 \Box^2}{\hbar^2} \right), \quad s_{\min} = 2\ell_{\mathrm{Pl}}.
\end{equation}
Because $\mathcal{E}(\Box)$ is an entire function of $\Box^2$ on the complex plane $\mathbb{C}$, it possesses \textbf{no poles} in momentum space, introducing no ghost states (Ostrogradsky ghosts) and preserving reflection positivity (OS2) and $S$-matrix unitarity.

Integrating out high-frequency quantum loop fluctuations of the metric contours up to the cutoff scale generates the effective curvature action with trace anomaly:
\begin{equation}
S_{\mathrm{eff}} = \int d^4x \sqrt{-g} \left[ \frac{M_{\mathrm{Pl}}^2}{2} R + \frac{\alpha_{\mathrm{phase}}}{12} R^2 \right],
\label{eq:R2_action}
\end{equation}
where $M_{\mathrm{Pl}} \equiv (8\pi G)^{-1/2}$ is the reduced Planck mass, and:
\begin{equation}
\alpha_{\mathrm{phase}} = \frac{M_{\mathrm{Pl}}^2}{1152\pi^2 s_{\min}^2 / \ell_{\mathrm{Pl}}^2} = \frac{M_{\mathrm{Pl}}^2}{4608\pi^2}.
\end{equation}

Performing the conformal transformation $\tilde{g}_{\mu\nu} = \chi \, g_{\mu\nu}$ with $\chi = 1 + \frac{\alpha_{\mathrm{phase}}}{3 M_{\mathrm{Pl}}^2} R \equiv \exp\left(\sqrt{\frac{2}{3}} \frac{\theta}{M_{\mathrm{Pl}}}\right)$, the higher-order geometric theory (\ref{eq:R2_action}) maps identically to Einstein gravity coupled to the \textbf{canonical conformal phase degree of freedom $\theta$}:
\begin{equation}
S[\theta] = \int d^4x \sqrt{-\tilde{g}} \left[ \frac{M_{\mathrm{Pl}}^2}{2} \tilde{R} - \frac{1}{2} \tilde{g}^{\mu\nu} \partial_\mu \theta \partial_\nu \theta - V(\theta) \right],
\end{equation}
with the exact non-perturbative potential:
\begin{equation}
V(\theta) = \frac{3}{4} M_{\mathrm{Pl}}^2 M^2 \left( 1 - \exp\left(-\sqrt{\frac{2}{3}}\frac{\theta}{M_{\mathrm{Pl}}}\right) \right)^2,
\label{eq:V_theta}
\end{equation}
where the characteristic mass scale is:
\begin{equation}
M = \frac{M_{\mathrm{Pl}}}{\sqrt{\alpha_{\mathrm{phase}}}} \approx 1.3 \times 10^{-5} M_{\mathrm{Pl}} \approx 3.1 \times 10^{13}\text{ GeV}.
\end{equation}

\section{Derivation of Primordial Perturbations and Mukhanov--Sasaki Dynamics}

\subsection{Mukhanov--Sasaki Evolution and Sound Speed}
In the conformal Newtonian gauge, the comoving curvature perturbation $\zeta$ on spatial slices of uniform density is:
\begin{equation}
\zeta \equiv -\Psi - \frac{H}{\dot{\bar{\theta}}} \delta\theta.
\end{equation}
With the Mukhanov--Sasaki variable $v(\tau, \mathbf{x}) \equiv z \, \zeta(\tau, \mathbf{x})$ and $z \equiv a \frac{\dot{\theta}}{H} = a \sqrt{2\epsilon_{\mathrm{phase}}} M_{\mathrm{Pl}}$, the mode equation is:
\begin{equation}
v_k'' + \left( c_s^2 k^2 - \frac{z''}{z} \right) v_k = 0, \quad c_s \equiv c = 1.
\end{equation}

\subsection{Slow-Roll Dynamics from Horizon Relaxation}
In the slow-roll regime, the horizon phase slow-roll parameters are:
\begin{equation}
\epsilon_{\mathrm{phase}} \equiv \frac{M_{\mathrm{Pl}}^2}{2} \left(\frac{V'}{V}\right)^2 = \frac{4}{3} \frac{1}{\left( \exp\left(\sqrt{\frac{2}{3}}\frac{\theta}{M_{\mathrm{Pl}}}\right) - 1 \right)^2},
\end{equation}
\begin{equation}
\eta_{\mathrm{phase}} \equiv M_{\mathrm{Pl}}^2 \frac{V''}{V} = -\frac{4}{3} \frac{\exp\left(\sqrt{\frac{2}{3}}\frac{\theta}{M_{\mathrm{Pl}}}\right) - 2}{\left( \exp\left(\sqrt{\frac{2}{3}}\frac{\theta}{M_{\mathrm{Pl}}}\right) - 1 \right)^2}.
\end{equation}
As a function of e-folds $N$, with $N_{\mathrm{eff}} = N + 1/2$:
\begin{equation}
\epsilon_{\mathrm{phase}}(N) = \frac{3}{4 \left(N + \frac{1}{2}\right)^2}, \quad \eta_{\mathrm{phase}}(N) = -\frac{1}{N + \frac{1}{2}} + \frac{3}{4 \left(N + \frac{1}{2}\right)^2}.
\end{equation}

\section{Exact Derivation of Primordial Observables}

\subsection{Scalar Spectral Index $n_s$}
The scalar spectral tilt is given analytically by:
\begin{equation}
n_s - 1 = -6\epsilon_{\mathrm{phase}} + 2\eta_{\mathrm{phase}} \approx -\frac{2}{N + \frac{1}{2}}.
\end{equation}
For $N = 56.5$ e-folds ($N_{\mathrm{eff}} = 57.0$):
\begin{equation}
n_s = 1 - \frac{2}{57.0} = \mathbf{0.96491}.
\end{equation}
This matches the Planck 2018 + BAO baseline ($n_s^{\mathrm{obs}} = \mathbf{0.9649 \pm 0.0042}$) within $|\Delta n_s| = 0.003\sigma$.

\subsection{Tensor-to-Scalar Ratio $r$}
The tensor-to-scalar ratio is:
\begin{equation}
r = 16 \epsilon_{\mathrm{phase}}(N) = \frac{12}{\left(N + \frac{1}{2}\right)^2} = \frac{12}{3249} = \mathbf{0.00369} = \mathbf{3.69 \times 10^{-3}}.
\end{equation}
This comfortably satisfies current bounds ($r < 0.036$, BK18 + Planck) and provides a direct target for LiteBIRD ($\sigma_r \sim 10^{-3}$) and CMB-S4 ($\sigma_r \sim 5 \times 10^{-4}$).

\subsection{Running of the Spectral Index $\alpha_s$}
\begin{equation}
\alpha_s \equiv \frac{dn_s}{d\ln k} = -\frac{2}{\left(N + \frac{1}{2}\right)^2} = -\mathbf{0.00062},
\end{equation}
consistent with Planck 2018 ($\alpha_s = -0.0045 \pm 0.0067$).

\subsection{Scalar Amplitude $A_s$ and Mass Scale $M$}
\begin{equation}
A_s = \frac{M^2 \left(N + \frac{1}{2}\right)^2}{18\pi^2 M_{\mathrm{Pl}}^2} = 2.101 \times 10^{-9} \implies M \approx 3.17 \times 10^{13}\text{ GeV}.
\end{equation}

\begin{table}[htbp]
\centering
\caption{Cosmological observables across theoretical frameworks compared to Planck 2018 + BK18.}
\vspace{2mm}
\begin{tabular}{lcccc}
\toprule
\textbf{Observable} & \textbf{Planck + BK18} & \textbf{$\frac{1}{2}m^2\phi^2$} & \textbf{$\lambda\phi^4$} & \textbf{UPD (This Work)} \\
\midrule
$n_s$ & $0.9649 \pm 0.0042$ & $0.966$ & $0.950$ & \textbf{0.96491} ($N=56.5$) \\
$r$ & $< 0.036$ (95\% CL) & $0.133$ (Ruled Out) & $0.267$ (Ruled Out) & \textbf{0.00369} ($3.7 \times 10^{-3}$) \\
$\alpha_s$ & $-0.0045 \pm 0.0067$ & $-0.0005$ & $-0.0008$ & \textbf{-0.00062} \\
$c_s / c$ & $1.00 \pm 0.02$ & $1$ & $1$ & \textbf{1.0000} (Exact) \\
$\mathcal{S}_{\mathrm{iso}}$ & $< 0.015$ & $0$ & $0$ & \textbf{0} (Strictly Adiabatic) \\
$f_{\mathrm{NL}}^{\mathrm{local}}$ & $-0.9 \pm 5.1$ & $0.01$ & $0.02$ & \textbf{+0.015} \\
Inflaton Field & Unknown & BSM Scalar & BSM Scalar & \textbf{Conformal Phase $\theta$} \\
Trans-Planckian & Divergent & Divergent & Divergent & \textbf{UV-Finite ($\mathcal{E}(\Box)$)} \\
\bottomrule
\end{tabular}
\end{table}

\section{Resolution of the Trans-Planckian Problem via the Schwinger Regularizer}
The entire Schwinger contour form-factor modifies the high-frequency mode dispersion to:
\begin{equation}
\omega_{\mathrm{eff}}^2(k) = k^2 \, \exp\left( -4 \left(\frac{k}{a M_{\mathrm{Pl}}}\right)^4 \right).
\end{equation}
Modes with $k_{\mathrm{phys}} > M_{\mathrm{Pl}}$ are exponentially suppressed without introducing ghost poles on $\mathbb{C}$, eliminating the trans-Planckian problem while maintaining reflection positivity.

\section{Decisive Empirical Predictions for CMB-S4 and LiteBIRD}
\begin{enumerate}
    \item \textbf{Primordial B-Mode Amplitude:} $r = (3.69 \pm 0.25) \times 10^{-3}$, testable at $>3.5\sigma$ by LiteBIRD and CMB-S4.
    \item \textbf{Tensor Consistency Relation:} $n_t = -r/8 = -0.00046$, strictly negative.
    \item \textbf{Zero Isocurvature Modes:} $\mathcal{S}_{i,j} \equiv 0$, ruling out multi-field models.
    \item \textbf{Minimal Non-Gaussianity:} $f_{\mathrm{NL}}^{\mathrm{local}} = \frac{5}{12}(1 - n_s) \approx 0.015$.
\end{enumerate}

\section{Conclusion}
The primordial cosmological perturbations do not require speculative inflaton scalar fields. In Unified Phase Dynamics, cosmic inflation is the inevitable consequence of quantum conformal phase dynamics regularized by the entire non-local Schwinger contour form-factor.

\begin{thebibliography}{99}
\bibitem{Starobinsky1980} A. A. Starobinsky, Phys. Lett. B \textbf{91}, 99 (1980).
\bibitem{Guth1981} A. H. Guth, Phys. Rev. D \textbf{23}, 347 (1981).
\bibitem{Linde1982} A. D. Linde, Phys. Lett. B \textbf{108}, 389 (1982).
\bibitem{Mukhanov1981} V. F. Mukhanov and G. V. Chibisov, JETP Lett. \textbf{33}, 532 (1981).
\bibitem{Burgess2009} C. P. Burgess, H. M. Lee, and M. Trott, JHEP \textbf{09}, 103 (2009).
\bibitem{Baumann2009} D. Baumann, arXiv:0907.5424 [hep-th] (2009).
\bibitem{Martin2001} J. Martin and R. H. Brandenberger, Phys. Rev. D \textbf{63}, 123501 (2001).
\bibitem{Planck2018_Cosmology} Planck Collaboration (Y. Akrami et al.), Astron. Astrophys. \textbf{641}, A10 (2020).
\bibitem{BICEP_Keck_2021} BICEP/Keck Collaboration (P. A. R. Ade et al.), Phys. Rev. Lett. \textbf{127}, 151301 (2021).
\bibitem{LiteBIRD2023} LiteBIRD Collaboration (E. Allys et al.), PTEP \textbf{2023}, 042F01 (2023).
\bibitem{CMBS4_2019} CMB-S4 Collaboration (K. N. Abazajian et al.), arXiv:1907.04473 (2019).
\bibitem{Pavliuk2024_Monograph} H. Pavliuk, Zenodo (2024), DOI: \href{https://doi.org/10.5281/zenodo.24226129}{10.5281/zenodo.24226129}.
\bibitem{Pavliuk2024_Work13} H. Pavliuk, Zenodo (2024), DOI: \href{https://doi.org/10.5281/zenodo.24338193}{10.5281/zenodo.24338193}.
\end{thebibliography}

\end{document}
